\documentclass[11pt,letterpaper]{article}

\usepackage[T1]{fontenc}
\usepackage[utf8]{inputenc}
\usepackage[margin=0.85in]{geometry}
\usepackage{graphicx}
\usepackage[table]{xcolor}
\usepackage{array}
\usepackage{booktabs}
\usepackage{longtable}
\usepackage{tabularx}
\usepackage{enumitem}
\usepackage{caption}
\usepackage{fancyhdr}
\usepackage{microtype}
\usepackage{hyperref}

\definecolor{SharedNavy}{HTML}{17324D}
\definecolor{SharedTeal}{HTML}{087E96}
\definecolor{SharedPale}{HTML}{EDF3F7}

\hypersetup{
  colorlinks=true,
  linkcolor=SharedNavy,
  urlcolor=SharedTeal,
  citecolor=SharedNavy,
  pdftitle={SharedHome Milestone 2 Design and Prototype Report},
  pdfauthor={Ramz, Hamza, and Karam}
}

\pagestyle{fancy}
\fancyhf{}
\fancyhead[R]{\footnotesize\textcolor{SharedNavy}{SHAREDHOME \textbar\ MILESTONE 2}}
\fancyfoot[C]{\footnotesize\thepage}
\renewcommand{\headrulewidth}{0pt}
\renewcommand{\footrulewidth}{0pt}

\setlength{\parindent}{0pt}
\setlength{\parskip}{0.55em}
\setlength{\tabcolsep}{4.5pt}
\renewcommand{\arraystretch}{1.22}
\setlist[itemize]{leftmargin=1.35em,itemsep=0.2em,topsep=0.25em}
\captionsetup{font=small,labelfont=bf}

\newcolumntype{Y}{>{\raggedright\arraybackslash}X}
\newcolumntype{P}[1]{>{\raggedright\arraybackslash}p{#1}}
\newcommand{\thead}[1]{\cellcolor{SharedNavy}\color{white}\bfseries #1}
\newcommand{\reportimage}[2]{%
  \IfFileExists{#1}{\includegraphics[width=#2]{#1}}{%
    \fbox{\parbox[c][2.2in][c]{#2}{\centering Missing image\\\texttt{#1}}}}}

\begin{document}

\begin{titlepage}
  \vspace*{0.7in}
  {\Huge\bfseries SharedHome Milestone 2 Design and Prototype Report\par}
  \vspace{0.35in}
  {\Large\color{SharedNavy}Role adaptive smart home interaction for administrators older adults and children\par}

  \vspace{0.55in}
  \begin{tabular}{@{}ll@{}}
    \textbf{Course} & Human Computer Interaction\\[0.35em]
    \textbf{Term} & Fall 2026\\[0.35em]
    \textbf{Milestone} & Solution Design and Prototype Development\\[0.35em]
    \textbf{Team} & Ramz, Hamza, Karam
  \end{tabular}

  \vfill
  This report documents the evidence chain, conceptual design, prototype development,
  and iterative refinement completed for SharedHome through the Iteration 3c candidate.
  \vspace{0.7in}
\end{titlepage}

\section{Project Summary}

SharedHome is a mobile smart-home prototype for multigenerational households. It changes
both interface complexity and action authority while keeping adaptations and restrictions
visible. The current Iteration 3c candidate combines user-selected Basic, Guided, and Full
experiences with explainable restrictions, scoped access requests, shared temperature
negotiation, and plain-language home rules.

Two formative reviews support continuing with role-aware layouts. Participants initially
felt more comfortable but still needed time to adjust. After brief role-specific guidance,
stable control placement, and clearer terminology were added, participants showed faster
learning, adaptation, and completion. They responded positively to icons paired with text
labels, and one participant described the prototype as a "convenient experience." These
are qualitative findings because numerical task times, error counts, and participant-level
outcomes were not recorded. Iteration 3c has not yet been evaluated and is reported as a
design hypothesis.

\section{Research Basis}

\subsection{Initial Participants}

\rowcolors{2}{white}{SharedPale}
\begin{tabularx}{\textwidth}{P{0.85in} P{1.45in} Y P{1.05in}}
  \thead{Participant} & \thead{Role} & \thead{Recorded finding} & \thead{Evidence}\\
  Ramz & Technically comfortable administrator & Reported no difficulty using IoT controls & Direct interview\\
  Murium & Older adult, age 60 & Has difficulty navigating smart-home controls & Direct interview\\
  Ayun & Child, age 10 & Requires limits on available controls & Direct interview\\
\end{tabularx}

\subsection{Evidence Boundaries}

The three interviews establish differences in technical confidence and the need for child
restrictions. They do not independently establish that users need restriction explanations,
in-app access requests, temporary approvals, or progressive disclosure. Those interactions
began as design hypotheses and were explored through the prototype. Wi-Fi and front-door
controls are also hypothesis devices rather than devices reported in the initial interviews.

\begin{itemize}
  \item The exploratory sample is small and drawn from one household context.
  \item Ramz is both a team member and participant, which may introduce response bias.
  \item The interview date, exact questions, and complete quotations are not recorded.
  \item The study record should document guardian permission and age-appropriate assent for Ayun.
\end{itemize}

\section{Design Requirements}

The requirements follow the course chain from user problem to research finding, design
requirement, HCI principle, design decision, and prototype feature. The evidence status
distinguishes interview findings from baseline usability requirements and hypotheses.

\small
\rowcolors{2}{white}{SharedPale}
\begin{longtable}{P{0.45in} P{1.05in} P{1.7in} P{1.35in} P{1.6in}}
  \thead{ID} & \thead{Evidence} & \thead{Requirement} & \thead{HCI rationale} & \thead{Prototype response}\\
  \endhead
  DR1 & Interview supported & Present role-appropriate interface complexity & Flexibility and cognitive-load reduction & Separate Administrator, Elder, and Child views\\
  DR2 & Interview supported & Limit Murium's view to essential actions & Hick's law and minimalist design & Large room-scoped device controls\\
  DR3 & Accessibility baseline & Make primary actions clear and easy to select & Visibility, consistency, and target size & Large labeled controls with visible state\\
  DR4 & Interview supported & Restrict selected child actions & Error prevention and safe exploration & Visible locked Wi-Fi, AC, and front-door controls\\
  DR5 & Design hypothesis & Explain every unavailable action & System status and problem diagnosis & Plain-language locked-control explanation\\
  DR6 & Design hypothesis & Let a restricted user request access & User control and feedback & Request action with pending and resolved status\\
  DR7 & Design hypothesis & Let the Administrator choose approval scope & Error prevention and reversibility & Once, one hour, always, and deny options\\
  DR8 & Usability baseline & Confirm each device or permission action & Visibility of system status & State changes, badges, and feedback messages\\
  DR9 & Iteration feedback & Support gradual learning and stable use & Consistency and recognition & Stable layout with dismissible guidance\\
  DR10 & Iteration feedback & Pair device symbols with explicit labels & Recognition over recall & Icon, name, and state shown together\\
\end{longtable}
\normalsize

\section{Conceptual Design}

\subsection{Personas}

\textbf{Murium.} Murium is 60 years old and has low confidence with technology. She uses
common household devices but finds smart-home navigation difficult. Her interface emphasizes
a few room-level controls, large targets, plain labels, and optional access to additional
functions.

\textbf{Ayun.} Ayun is 10 years old and needs limits on available controls. His interface
keeps allowed devices easy to reach and presents restricted controls as visible locked items.
Wi-Fi and the front door are included as hypothesis devices for testing whether permission
boundaries are understood.

\textbf{Administrator.} The Administrator is comfortable with smart-home controls and needs
full device access, permission management, and efficient request handling. The interface
prioritizes pending requests and shows the scope of each approval.

\subsection{Scenarios and Task Flows}

\rowcolors{2}{white}{SharedPale}
\begin{tabularx}{\textwidth}{P{1.45in} Y P{1.75in}}
  \thead{Scenario} & \thead{Main flow} & \thead{Success condition}\\
  Murium controls comfort & Open Elder view, select Fan, turn it on, and confirm state & Complete without help or entering settings\\
  Ayun requests Wi-Fi & Open Child view, select locked Wi-Fi, read the explanation, and request access & Explain the restriction and send one request\\
  Administrator resolves request & Open Administrator view, review the request, select one hour, and confirm & Grant only the intended device and duration\\
\end{tabularx}

\subsection{Information Architecture}

\begin{itemize}
  \item Prototype role switcher.
  \item Administrator dashboard with requests, devices, access settings, and shared decisions.
  \item Murium dashboard with stable essential controls and progressive disclosure.
  \item Ayun dashboard with allowed devices and visible locked controls.
  \item Locked-control explanation, request action, request status, and approval scope.
\end{itemize}

\section{HCI Principles and Design Decisions}

\small
\rowcolors{2}{white}{SharedPale}
\begin{tabularx}{\textwidth}{P{1.3in} P{2.25in} Y}
  \thead{Principle} & \thead{Design decision} & \thead{Prototype implementation}\\
  Hick's law & Reduce visible choices in the simpler experience & Three primary controls with additional controls disclosed later\\
  Cognitive-load reduction & Keep each role focused on likely tasks & Role-specific content and short contextual guidance\\
  Recognition over recall & Pair familiar symbols with text & Icon, device name, and state appear together\\
  Visibility of system status & Show the result of device and request actions & State badges, request counts, and feedback messages\\
  Error prevention & Stop restricted or high-consequence actions before execution & Locked controls and scoped Administrator approval\\
  User control and freedom & Let users dismiss guidance and reverse safe changes & Close, Escape, Not now, and Undo actions\\
  Progressive disclosure & Delay secondary controls until requested & Experience modes and additional-control sections\\
  Accessible interaction & Provide generous targets and keyboard behavior & Large tiles, focus transfer, and Escape dismissal\\
\end{tabularx}
\normalsize

\section{Interactive Prototype}

SharedHome is implemented as a sequence of clickable web prototypes with a role switcher.
Cross-role state lets a request made by Ayun appear in the Administrator view and an approval
return to Ayun's view. Iteration 3c merges interface adaptation from Iteration 3a with
household coordination from Iteration 3b while preserving the phone-based format.

\textbf{Device scope.} Interview-grounded devices include lights, air conditioning,
television, and fans. Wi-Fi and the front door are hypothesis devices used to test permission
explanations, safety boundaries, and approval scope.

\begin{figure}[htbp]
  \centering
  \begin{minipage}[t]{0.47\textwidth}
    \centering
    \reportimage{report_assets/iteration-2-elder.png}{0.88\linewidth}\\
    \small Elder view
  \end{minipage}\hfill
  \begin{minipage}[t]{0.47\textwidth}
    \centering
    \reportimage{report_assets/iteration-2-child.png}{0.88\linewidth}\\
    \small Child view
  \end{minipage}
  \caption{Iteration 2 role-specific Elder and Child layouts}
\end{figure}

\begin{figure}[htbp]
  \centering
  \reportimage{report_assets/iteration-2-admin.png}{0.43\textwidth}
  \caption{Iteration 2 Administrator dashboard and access management}
\end{figure}

\clearpage
\section{Iterative Refinement}

The prototype developed through a baseline concept, two feedback-led refinements, and three
Iteration 3 branches. The table records what changed, why the team made the change, what it
achieved, and the limitation that motivated the next step.

\small
\rowcolors{2}{white}{SharedPale}
\begin{longtable}{P{0.55in} P{2.65in} P{3.15in}}
  \thead{Iteration} & \thead{Change and purpose} & \thead{Result and remaining limitation}\\
  \endhead
  0 & Established Administrator, Elder, and Child views with common device controls, visible locked items, and a request and approval flow. & Made the role-aware concept testable. Views were fixed, guidance was thin, and learning and adaptation were not addressed.\\
  1 & Tested the role layouts with the three participants and retained Wi-Fi and the front door as permission hypotheses. & Participants felt more comfortable but still needed time to adapt. No task times, errors, or attributed quotations were recorded.\\
  2 & Added short guidance, stable device placement, consistent terminology, icon-plus-label controls, dialog focus, and Escape dismissal. & Participants reported faster learning and completion and a convenient experience. Evidence remained qualitative.\\
  3a & Introduced Basic, Guided, and Full experiences, transparent recommendations, four-part restriction explanations, and reversible control. & Strengthened adaptation, accessibility, and autonomy but did not address shared decisions or Administrator workload.\\
  3b & Added repeated-request rule suggestions, living-room AC preference negotiation, and plain-language home rules. & Strengthened coordination and fairness questions but omitted several features from 3a.\\
  3c & Merged 3a and 3b and corrected schedule labeling and rule-suggestion safety conditions. & Current evaluation candidate. Adaptation, PACE recommendations, and schedules are simulated and unvalidated.\\
\end{longtable}
\normalsize

\subsection{Participant Feedback Through Iteration 2}

The first review suggested that role-specific layouts reduced some discomfort, but users
still needed time to learn the arrangement. Iteration 2 kept essential controls in stable
positions and placed short guidance near unfamiliar actions. The follow-up feedback indicated
faster learning, adaptation, and completion. Because the team did not record task-level
measurements or the speaker of the quotation, these observations support continued design
work but do not establish a measured improvement.

\begin{center}
  \large\bfseries "Convenient experience"
\end{center}

\section{Iteration 3c Candidate Prototype}

Iteration 3c combines the strongest parts of both Iteration 3 branches. It keeps low-risk
essential controls stable while allowing a user to select the amount of detail and guidance.
It also makes household rules and shared decisions visible across roles.

\begin{itemize}
  \item Murium can select Basic, Guided, or Full without changing her authority. A simulated recommendation explains why more controls are suggested and requires approval.
  \item Ayun sees locked Wi-Fi, air-conditioning, and front-door items. The explanation identifies the reason, decision owner, available next action, and access duration.
  \item The Administrator sees access requests and shared AC decisions as separate workloads and can review limited approval options.
  \item Murium and Ayun can state different AC preferences, while the Administrator can select a compromise or one person's preference.
  \item Home rules show who set each restriction and why. In Murium's view they remain behind additional controls.
\end{itemize}

\subsection{Safety and Validity Corrections}

\begin{itemize}
  \item Recurring suggestions count only approved requests. Pending or denied requests do not count, and no suggestion appears while matching access is active.
  \item Accepting a recurring rule resolves a matching pending request. The front door is excluded from recurring suggestions.
  \item The prototype labels schedules, familiarity, and PACE recommendations as simulated. It does not claim real-time scheduling or machine learning.
\end{itemize}

Ability-based adaptation, a household co-pilot, multimodal language support, and a
vendor-neutral device layer remain future work rather than implemented features.

\begin{figure}[htbp]
  \centering
  \begin{minipage}[t]{0.47\textwidth}
    \centering
    \reportimage{report_assets/iteration-3c-elder.png}{0.88\linewidth}\\
    \small Murium's Guided experience and transparent recommendation
  \end{minipage}\hfill
  \begin{minipage}[t]{0.47\textwidth}
    \centering
    \reportimage{report_assets/iteration-3c-child.png}{0.88\linewidth}\\
    \small Ayun's allowed and locked controls with shared AC context
  \end{minipage}
  \caption{Iteration 3c role-specific views}
\end{figure}

\begin{figure}[htbp]
  \centering
  \reportimage{report_assets/iteration-3c-admin.png}{0.43\textwidth}
  \caption{Iteration 3c Administrator requests, devices, and shared decisions}
\end{figure}

\clearpage
\section{Project Contribution and Novelty}

Role permissions, child restrictions, device icons, and temporary access are not individually
novel. Commercial systems already use owner, member, resident, or guest permission models
\cite{googlehome,applehome}. SharedHome's proposed contribution is the combined interaction
model: the system changes presentation and permitted action, explains each boundary, supports
a response inside the interface, and keeps recommendations under user control. This is a
course-project contribution and research hypothesis, not a claim that no existing product
contains any individual feature.

\small
\rowcolors{2}{white}{SharedPale}
\begin{tabularx}{\textwidth}{P{1.1in} P{1.55in} Y P{1.45in}}
  \thead{Dimension} & \thead{Common baseline} & \thead{SharedHome Iteration 3c} & \thead{Evidence status}\\
  Interface complexity & A role mainly changes access rights. & The user selects Basic, Guided, or Full while essential controls stay in place. & Grounded in Murium's difficulty; 3c is untested.\\
  Restriction feedback & A control may be unavailable or limited. & The interface explains why, who decides, what can be requested, and when access expires. & Design hypothesis.\\
  Adaptation & Settings may be personalized manually. & A recommendation gives its reason and requires Accept, Not now, or Undo. & Simulated hypothesis.\\
  Household conflict & An owner controls the device. & Different preferences remain visible and an Administrator resolves the AC decision. & Fairness hypothesis.\\
  Repeated requests & The owner handles requests individually. & Two approved Wi-Fi or AC requests can trigger a simulated recurring rule. & Interruption hypothesis.\\
\end{tabularx}
\normalsize

\subsection{PACE Design Model}

PACE is the team's proposed description of transparent adaptation. It is not a validated
general model. Each dimension changes presentation, guidance, approval scope, or feedback.
The approach also responds to older-adult smart-home research that emphasizes customization
and avoiding age-based assumptions \cite{ghorayeb}.

\rowcolors{2}{white}{SharedPale}
\begin{tabularx}{\textwidth}{P{0.95in} P{2.25in} Y}
  \thead{Dimension} & \thead{Design question} & \thead{Prototype example}\\
  Person & What are this person's preferences, accessibility needs, and authority? & Murium selects Guided mode while Ayun retains limited authority.\\
  Action risk & What could happen if the action is performed incorrectly? & Lights are low risk, Wi-Fi is moderate risk, and the front door is excluded from recurring suggestions.\\
  Context & When and for how long is the action needed? & Ayun requests one hour of Wi-Fi; the recurring 4--6 p.m. schedule is visibly simulated.\\
  Experience & How familiar is the person with the task? & A simulated recommendation appears after successful use and requires explicit acceptance.\\
\end{tabularx}

\section{Phase 2 Deliverable Status}

The current work covers every Phase 2 deliverable in the course brief. Items marked for
validation belong to the Phase 3 study and are not evaluated outcomes.

\small
\rowcolors{2}{white}{SharedPale}
\begin{tabularx}{\textwidth}{P{1.35in} Y P{2.05in}}
  \thead{Course requirement} & \thead{Completed work} & \thead{Status and next evidence}\\
  Design requirements & Ten requirements connect participant evidence or hypotheses to HCI rationale and prototype responses. & Completed for Phase 2; follow-up evidence is needed for restriction and request hypotheses.\\
  Conceptual design & Personas, scenarios, task flows, information architecture, and role-specific wireframes are documented. & Completed; the Elder label evolved into user-selected experience levels.\\
  HCI principles & Visibility, feedback, consistency, error prevention, recognition, cognitive-load reduction, progressive disclosure, user control, and accessibility are traced to features. & Completed and traceable.\\
  Prototype development & Clickable Iterations 0, 1, 2, 3a, 3b, and 3c support role switching and cross-role state. & Iteration 3c is the current candidate.\\
  Iterative refinement & The log records changes, reasons, feedback, unresolved limits, and the 3a and 3b merge. & Completed through 3c; formal evaluation remains pending.\\
\end{tabularx}
\normalsize

\section{Evaluation Readiness}

Phase 3 should compare Iteration 3c with a stripped-down baseline containing the same devices
but a single-owner model, no restriction explanations, no adaptation recommendations, and no
shared negotiation. The baseline has not yet been built. A within-participant comparison can
reveal usability problems, but the three-person sample cannot support broad population claims.
Task order should be counterbalanced where possible.

\subsection{Research Questions}

\begin{itemize}
  \item Does transparent adaptation based on person, action risk, context, and experience improve task performance and permission understanding without reducing autonomy?
  \item Do explainable restrictions help Ayun understand why an action is blocked, who manages it, what he can request, and when access expires?
  \item Do rule suggestions reduce Administrator interruptions, and does shared AC negotiation improve perceived fairness?
\end{itemize}

\subsection{Planned Tasks and Measures}

\small
\rowcolors{2}{white}{SharedPale}
\begin{tabularx}{\textwidth}{P{2.05in} P{2.1in} Y}
  \thead{Persona and task} & \thead{Performance evidence} & \thead{Experience evidence}\\
  Murium selects Guided mode, turns on the fan, responds to a shortcut suggestion, and reverses it. & Success, time, errors, assistance, hesitation, and correct use of Accept, Not now, and Undo. & Ease, confidence, perceived control, and whether the suggestion felt helpful or patronizing.\\
  Ayun opens locked Wi-Fi, explains the restriction, requests one hour, and identifies the expiry. & Correct answers for why, who, next action, and duration; success, time, errors, and assistance. & Clarity, fairness, and whether the restriction felt like a dead end.\\
  Ramz reviews Wi-Fi, handles the shared AC decision, and evaluates a recurring-rule suggestion. & Correct person, device, action, duration, and risk; completion time; requests handled; interruptions. & Workload, confidence, recommendation usefulness, and perceived fairness.\\
\end{tabularx}
\normalsize

For every task, record the date, participant, condition, completion result, time in seconds,
errors, assistance, hesitation points, chosen mode or permission, and one attributed
quotation. Add brief ratings for clarity, autonomy, and fairness. If the course permits,
report a standardized usability questionnaire separately from qualitative comments.

\section{Limitations and Ethics}

\begin{itemize}
  \item The evidence comes from three participants in one household context, so results will be exploratory.
  \item Ramz is both a team member and participant, which may introduce response bias.
  \item The interview date, exact questions, recruitment details, and complete quotations remain undocumented.
  \item The study record should document guardian permission and age-appropriate assent for Ayun.
  \item Earlier reviews did not record numerical task times, error counts, or participant-level results, and the speaker of "convenient experience" is unknown.
  \item Wi-Fi and the front door are test scenarios rather than devices reported in the original interviews.
  \item PACE, adaptation, scheduling, and recurring-rule behavior are simulated. No machine learning or real-time schedule is operating.
\end{itemize}

\section{Team Responsibilities}

\rowcolors{2}{white}{SharedPale}
\begin{tabularx}{\textwidth}{P{1.2in} Y}
  \thead{Owner} & \thead{Deliverables}\\
  Ramz & Design requirements, evidence-to-feature chain, and HCI principle justification.\\
  Hamza & Wireframes, information architecture, interactive prototypes, and prototype implementation.\\
  Karam & Personas, scenarios, user flows, task models, evaluation records, and iteration log.\\
\end{tabularx}

\section{Conclusion}

SharedHome now meets the Phase 2 goal of translating user evidence into an interactive,
theory-grounded prototype. Iterations 0 through 2 established role-aware layouts, stable
controls, contextual guidance, and icon-plus-label presentation. Iteration 3c adds transparent
user-controlled adaptation, explainable restrictions, scoped access, and household
negotiation. Its value has not yet been established through formal testing. The next step is
to build the stripped-down baseline and evaluate both conditions with recorded task
performance, permission comprehension, autonomy, Administrator interruptions, and fairness.

\begin{thebibliography}{9}

\bibitem{applehome}
Apple. (2026).
\newblock \emph{Allow others to control accessories in your home}.
\newblock Apple Support.
\newblock \url{https://support.apple.com/guide/iphone/iphcbaf7e8f3/ios}

\bibitem{ghorayeb}
Ghorayeb, A., Comber, R., and Gooberman-Hill, R. (2023).
\newblock Older adults' perspectives of smart home technology: Are we developing the technology that older people want?
\newblock \emph{International Journal of Human-Computer Studies}, 170, 102979.
\newblock \url{https://doi.org/10.1016/j.ijhcs.2022.102979}

\bibitem{googlehome}
Google. (2026).
\newblock \emph{Share a home and devices in the Google Home app}.
\newblock Google Nest Help.
\newblock \url{https://support.google.com/googlenest/answer/9155535}

\bibitem{hick}
Hick, W. E. (1952).
\newblock On the rate of gain of information.
\newblock \emph{Quarterly Journal of Experimental Psychology}, 4(1), 11--26.
\newblock \url{https://doi.org/10.1080/17470215208416600}

\bibitem{nielsen}
Nielsen, J. (1994).
\newblock \emph{10 usability heuristics for user interface design}.
\newblock Nielsen Norman Group.
\newblock \url{https://www.nngroup.com/articles/ten-usability-heuristics/}

\bibitem{norman}
Norman, D. A. (2013).
\newblock \emph{The design of everyday things}. Revised and expanded edition.
\newblock Basic Books.

\bibitem{wcag}
World Wide Web Consortium. (2023).
\newblock \emph{Web Content Accessibility Guidelines 2.2}.
\newblock \url{https://www.w3.org/TR/WCAG22/}

\end{thebibliography}

\end{document}
